\begin{table}[H]
\centering
\caption{Linearity Test: Connectivity and Firm Outcomes}
\label{tab:linearity_binstest}
\scriptsize
\begin{tabular}{llccccc}
\toprule\toprule
Outcome & Sample & $N$ & \shortstack{Bins\\(DPI)} & \shortstack{Sup-$t$\\stat} & $p$-value \\
\midrule
Log  wages & Year 2011 & $4{,}183$ & 28 & 0.577 & 0.985 \\
Log hourly wages & Year 2011 & $4{,}183$ & 28 & 1.406 & 0.498 \\
Log employment & Year 2011 & $4{,}195$ & 28 & 2.228 &  0.083 \\
\# CBA clauses & CBA period 1 & $4{,}356$ & 29 & 1.226 & 0.610 \\
\bottomrule
\end{tabular}
\begin{minipage}{\linewidth}
\footnotesize\vspace{4pt}
\textit{Notes:} This table reports Cattaneo, Crump, Farrell, and Feng (2024) sup-$t$ tests of the null hypothesis that the conditional expectation of each outcome is a degree-1 polynomial (linear) in connectivity. The cross-section is fixed at year 2011 (calendar-year outcomes) or CBA period 1 (number of CBA clauses). Controls included as covariates: industry, month-of-year, microregion, pre-treatment outcome quartile, pre-treatment employment quartile, and pre-treatment worker-flow quartile. Bins are set to 50; effective bins after dropping repeated quantile knots at $x=0$ are reported. Inference based on 2{,}000 simulation draws; grid of 50 evaluation points per bin. Standard errors are heteroskedasticity-robust.
\end{minipage}
\end{table}
